\chapter{Complexity — Path Walk Proof}
%%% generated by the blueprint pipeline from TCSlib.Complexity.SpaceComplexity.Logspace.PathWalkProof
\section{Overview}

Declarations of the Lean module \texttt{TCSlib.Complexity.SpaceComplexity.Logspace.PathWalkProof}.

\section{Declarations}

\begin{lemma}[Plain block comparison]
\lean{Complexity.LogProg.ps_tcmp}
\label{Complexity.LogProg.ps_tcmp}
\leanok
\uses{Complexity.LogProg.NSem.call_neg, Complexity.LogProg.NSem.call_pos, Complexity.LogProg.NSem.congr, Complexity.LogProg.NSem.half, Complexity.LogProg.NSem.inc, Complexity.LogProg.NSem.jodd_neg, Complexity.LogProg.NSem.jodd_pos, Complexity.LogProg.NSem.jz_neg, Complexity.LogProg.NSem.jz_pos, Complexity.LogProg.NSem.ret, Complexity.LogProg.PReg, Complexity.LogProg.PReg.Le, Complexity.LogProg.PReg.reg, Complexity.LogProg.PReg.v_0, Complexity.LogProg.PReg.v_6, Complexity.LogProg.PS, Complexity.LogProg.PS.of_eq, Complexity.LogProg.pathOr_bit, Complexity.LogProg.pathOr_len, Complexity.PlainAt, Complexity.bits_eq_cons, Complexity.plainAt_cons}
\difficulty{6}
Let $x$ be a finite bit string, let $B$ be a natural number, let $K$ be a proposition,
and let $\sigma=(p,n,a,b,s,t,tmp,st,nx)$ be a nine-register state. Suppose $2|x|+2\le
B$, every register of $\sigma$ is at most $B$, and $p\le |x|$. Assume that whenever the
suffix of $x$ beginning at position $p$ is exactly the binary word for $tmp$, in the
sense that $\operatorname{PlainAt}(x,p,\operatorname{bits}(tmp))$ holds, the program
started at label $w0$ with $p=|x|$ and $tmp=0$ terminates within the register bound and
has an accepting branch exactly when $K$ holds. Then the program started at label $ttt$
in state $\sigma$ terminates within the register bound and has an accepting branch
exactly when both $\operatorname{PlainAt}(x,p,\operatorname{bits}(tmp))$ and $K$ hold.
\end{lemma}

\begin{lemma}[Existential target phase]
\lean{Complexity.LogProg.ps_tphase}
\label{Complexity.LogProg.ps_tphase}
\leanok
\uses{Complexity.LogProg.NSem.congr, Complexity.LogProg.PReg, Complexity.LogProg.PReg.Le, Complexity.LogProg.PReg.reg, Complexity.LogProg.PReg.v, Complexity.LogProg.PReg.v_1, Complexity.LogProg.PReg.v_5, Complexity.LogProg.PS, Complexity.LogProg.PS.of_eq, Complexity.LogProg.nsem_copy, Complexity.LogProg.nsem_guess, Complexity.LogProg.pathA, Complexity.LogProg.pathOr, Complexity.LogProg.ps_tcmp, Complexity.PlainAt}
\difficulty{3}
Let $x$ be a finite bit string, let $B$ be a natural number, let $K$ be a predicate on
natural numbers, and let $\sigma=(p,n,a,b,s,t,tmp,st,nx)$ be a nine-register state.
Suppose $2|x|+2\le B$, every register of $\sigma$ is at most $B$, $t=0$, $tmp=0$, $p\le
|x|$, and $n\le |x|$. Assume that for every $j<n$, whenever
$\operatorname{PlainAt}(x,p,\operatorname{bits}(j))$ holds, the program started at label
$w0$ with target register $t=j$, input pointer $p=|x|$, and scratch register $tmp=0$ has
property $K(j)$. Then the program started at label $tg$ in state $\sigma$ terminates on
every branch with all registers at most $B$, and has an accepting branch exactly when
there is some $j<n$ such that $\operatorname{PlainAt}(x,p,\operatorname{bits}(j))$ and
$K(j)$ both hold.
\end{lemma}

\begin{lemma}[Doubled word scan and continuation]
\lean{Complexity.LogProg.ps_scmp}
\label{Complexity.LogProg.ps_scmp}
\leanok
\uses{Complexity.DblAt, Complexity.LogProg.NSem.call_neg, Complexity.LogProg.NSem.call_pos, Complexity.LogProg.NSem.congr, Complexity.LogProg.NSem.half, Complexity.LogProg.NSem.inc, Complexity.LogProg.NSem.jodd_neg, Complexity.LogProg.NSem.jodd_pos, Complexity.LogProg.NSem.jz_neg, Complexity.LogProg.NSem.jz_pos, Complexity.LogProg.NSem.ret, Complexity.LogProg.PReg, Complexity.LogProg.PReg.Le, Complexity.LogProg.PReg.reg, Complexity.LogProg.PReg.v_0, Complexity.LogProg.PReg.v_6, Complexity.LogProg.PS, Complexity.LogProg.PS.of_eq, Complexity.LogProg.pathOr_bit, Complexity.bits_eq_cons, Complexity.dblAt_cons}
\difficulty{6}
Let $x$ be a finite bit string, $B$ a natural number, $K$ a proposition, and
$σ=(p,n,a,b,s,t,tmp,st,nx)$ a program register state. Suppose every register of $σ$ is
at most $B$ and $p+2|\mathrm{bits}(tmp)|+2\le B$. Assume that whenever the word of
binary digits of $tmp$ is doubled at position $p$ in $x$, the program started at the
target stage with pointer $p+2|\mathrm{bits}(tmp)|+2$ and scratch register zero
terminates with all registers bounded by $B$ and has an accepting computation exactly
when $K$ holds. Then the program started at the scan stage in state $σ$ terminates with
all registers bounded by $B$ and has an accepting computation exactly when both that
doubled-word condition and $K$ hold.
\end{lemma}

\begin{lemma}[Existential phase search]
\lean{Complexity.LogProg.ps_sphase}
\label{Complexity.LogProg.ps_sphase}
\leanok
\uses{Complexity.DblAt, Complexity.LogProg.NSem.congr, Complexity.LogProg.PReg, Complexity.LogProg.PReg.Le, Complexity.LogProg.PReg.reg, Complexity.LogProg.PReg.v, Complexity.LogProg.PReg.v_1, Complexity.LogProg.PReg.v_4, Complexity.LogProg.PS, Complexity.LogProg.PS.of_eq, Complexity.LogProg.nsem_copy, Complexity.LogProg.nsem_guess, Complexity.LogProg.pathA, Complexity.LogProg.pathOr, Complexity.LogProg.ps_scmp, Turing.length_bits_le_self}
\difficulty{3}
Let $x$ be a finite bit string, $B$ a natural number, let $K(j)$ be a proposition for
each natural number $j$, and let $σ=(p,n,a,b,s,t,tmp,st,nx)$ be a program register
state. Suppose every register of $σ$ is at most $B$, $s=0$, $tmp=0$, and $p+2n+2\le B$.
Assume that for every $j<n$, if the binary word of $j$ is doubled at position $p$ in
$x$, then the program started at the target stage with source register $j$, pointer
$p+2|\mathrm{bits}(j)|+2$, and scratch register zero terminates with all registers
bounded by $B$ and has an accepting computation exactly when $K(j)$ holds. Then the
program started at the existential-search stage in state $σ$ terminates with all
registers bounded by $B$ and has an accepting computation exactly when there exists
$j<n$ for which the binary word of $j$ is doubled at position $p$ in $x$ and $K(j)$
holds.
\end{lemma}

\begin{lemma}[Matrix scan and suffix property]
\lean{Complexity.LogProg.ps_mat}
\label{Complexity.LogProg.ps_mat}
\leanok
\uses{Complexity.LogProg.NSem.call_neg, Complexity.LogProg.NSem.call_pos, Complexity.LogProg.NSem.clr, Complexity.LogProg.NSem.congr, Complexity.LogProg.NSem.inc, Complexity.LogProg.NSem.jeq_neg, Complexity.LogProg.NSem.jeq_pos, Complexity.LogProg.NSem.ret, Complexity.LogProg.PReg, Complexity.LogProg.PReg.Le, Complexity.LogProg.PReg.reg, Complexity.LogProg.PReg.v_0, Complexity.LogProg.PReg.v_1, Complexity.LogProg.PReg.v_2, Complexity.LogProg.PS, Complexity.LogProg.PS.of_eq, Complexity.LogProg.pathOr_bit, Complexity.LogProg.pathOr_len, Complexity.MatRest, Complexity.div_eq_self_iff, Complexity.divmod_succ, Complexity.lt_of_getD_true}
\difficulty{7}
Let $x$ be a finite Boolean word, $B$ a natural number, $K$ a proposition, $i$ a natural
number, and $\sigma$ a program register state whose nine register values are all at most
$B$. Suppose $2|x|+2\le B$, the vertex-count register $n$ satisfies $n\le |x|$, the
input pointer is $2n+2+2i$, the row register is $\lfloor i/n\rfloor$, the column
register is the remainder of $i$ upon division by $n$, and $i\le n^2$. If the bit of $x$
at position $2n+2+2n^2+1$, using false for an out-of-range position, is true, and the
path program started at label $sg$ with input pointer $2n+2+2n^2+2$, row register $n$,
and column register $0$ satisfies $PS$ with proposition $K$, then the program started at
label $m0$ with register state $\sigma$ satisfies $PS$ with proposition that $K$ holds
and, for every $j$ with $i\le j<n^2$, the bits at positions $2n+2+2j$ and $2n+2+2j+1$
agree (out-of-range bits count as false), while the bits at positions $2n+2+2n^2$ and
$2n+2+2n^2+1$ are false and true, respectively (again counting out-of-range bits as
false).
\end{lemma}

\begin{lemma}[Unary prefix phase]
\lean{Complexity.LogProg.ps_unary}
\label{Complexity.LogProg.ps_unary}
\leanok
\uses{Complexity.LogProg.NSem.call_neg, Complexity.LogProg.NSem.call_pos, Complexity.LogProg.NSem.congr, Complexity.LogProg.NSem.inc, Complexity.LogProg.NSem.ret, Complexity.LogProg.PReg, Complexity.LogProg.PReg.Le, Complexity.LogProg.PReg.reg, Complexity.LogProg.PReg.v_0, Complexity.LogProg.PS, Complexity.LogProg.PS.of_eq, Complexity.LogProg.pathOr_bit, Complexity.UnaryRest, Complexity.lt_of_getD_true}
\difficulty{6}
Let $x$ be a finite bit string, let $B$ be a natural number, let $K$ be a predicate on
natural numbers, and let $\sigma=(p,n,a,b,s,t,tmp,st,nx)$ be a nine-register state.
Suppose $2|x|+2\le B$, every register of $\sigma$ is at most $B$, $p=2n$, and $p\le
|x|$. Assume that for every $n'\ge n$, if $\operatorname{UnaryRest}(x,n,n')$ holds, then
the program started at label $m0$ with $p=2n'+2$ and vertex count $n'$ has property
$K(n')$. Then the program started at label $u0$ in state $\sigma$ terminates on every
branch with all registers at most $B$, and has an accepting branch exactly when there is
some $n'\ge n$ for which $\operatorname{UnaryRest}(x,n,n')$ and $K(n')$ both hold.
\end{lemma}

\begin{definition}[Zero register configuration]
\lean{Complexity.LogProg.PReg.zero}
\label{Complexity.LogProg.PReg.zero}
\leanok
\uses{Complexity.LogProg.PReg}
A register configuration consists of nine natural-number registers, and the zero
configuration is the one in which every register has value $0$.
\end{definition}

\begin{theorem}[Path program from the zero configuration]
\lean{Complexity.LogProg.ps_start}
\label{Complexity.LogProg.ps_start}
\leanok
\uses{Complexity.DblAt, Complexity.LogProg.NSem.congr, Complexity.LogProg.PReg.Le, Complexity.LogProg.PReg.zero, Complexity.LogProg.PS, Complexity.LogProg.ps_mat, Complexity.LogProg.ps_sphase, Complexity.LogProg.ps_tphase, Complexity.LogProg.ps_unary, Complexity.LogProg.ps_walk, Complexity.MatRest, Complexity.PathAcc, Complexity.PlainAt, Complexity.WalkLe, Complexity.lt_of_getD_true, Complexity.pathEdge, Turing.length_bits_le_self}
\difficulty{5}
For every finite bit string $x$, the fixed nondeterministic path program started with
all nine registers zero and register bound $4(|x|+1)$ terminates on every computation
branch without exceeding that bound, and has an accepting branch exactly when the
following condition holds. There is a natural number $n$ such that for each $i<n$, the
entries of $x$ at positions $2i$ and $2i+1$ are true, and the entries at $2n$ and $2n+1$
are respectively false and true; for every $j<n^2$, the entries at positions $2n+2+2j$
and $2n+3+2j$ agree, and the entries at positions $2n+2+2n^2$ and $2n+3+2n^2$ are
respectively false and true. There are also $s,t<n$ such that the binary word of $s$ is
doubled at position $2n+2+2n^2+2$ in $x$, the binary word of $t$ occurs starting at
position $2n+2+2n^2+2+2|\mathrm{bits}(s)|+2$ and ends exactly at the end of $x$, and
there is a walk from $s$ to $t$ of length at most $n$ whose vertices are less than $n$,
with each step from $u$ to $v$ satisfying $x_{2n+2+2(un+v)}=1$ (out-of-range entries are
false).
\end{theorem}

\begin{theorem}[The raw path language is in NL]
\lean{Complexity.pathRawLang_mem_NL}
\label{Complexity.pathRawLang_mem_NL}
\leanok
\uses{Complexity.LogProg.NHalts.mono, Complexity.LogProg.NSem.decides, Complexity.LogProg.PReg.v, Complexity.LogProg.PReg.zero, Complexity.LogProg.narm_decides_poly, Complexity.LogProg.pathA, Complexity.LogProg.pathA_npreS, Complexity.LogProg.pathAs, Complexity.LogProg.pathAs_mem, Complexity.LogProg.ps_start, Complexity.NL, Complexity.pathAcc_iff, Complexity.pathRawLang}
\difficulty{4}
The language of encodings of finite directed graphs with designated vertices, for which
the target is reachable from the source, belongs to nondeterministic logarithmic space.
\end{theorem}
